\section{Quadratic fitting and its transport law}\label{sec:fitting}

For the reference profile
\[
 q_{\lambda,\kappa}(u)=\lambda u+\frac{\kappa}{2}u^2,
 \qquad (\lambda,\kappa)\in\mathbb R^2,
\]
let \(H(\lambda,\kappa,r)\) be its midpoint data at base height zero
and width \(r>0\).  The entire-profile version of Lemma~\ref{lem:arc-data}
makes this definition valid for every such triple.  Set
\begin{equation}\label{eq:fit-notation}
 P=H_\lambda,\qquad Q=H_\kappa,\qquad
 R=P_r,\qquad S=Q_r,\qquad G=PS-QR.
\end{equation}
Unless stated otherwise, a subscript on a model function denotes an
ordinary partial derivative with its other arguments held fixed.
Smooth arc dependence gives smoothness of these functions for \(r>0\),
and Lemma~\ref{lem:arc-data} gives
\begin{equation}\label{eq:fit-open-data}
 |H|<r,\qquad |H_r|<1.
\end{equation}
The variations \(u\) and \(u^2/2\) of the profile are strictly positive
in the interior of every base-zero arc.  The positive denominator in
Lemma~\ref{lem:variation} therefore gives, already at this stage,
\begin{equation}\label{eq:fit-PQ-positive}
 P>0,\qquad Q>0.
\end{equation}

Reflecting the endpoint coordinate \(y\) negates the profile and the
midpoint, while preserving the width.  Consequently
\begin{equation}\label{eq:fit-reflection}
 H(-\lambda,-\kappa,r)=-H(\lambda,\kappa,r),\qquad
 H_r(-\lambda,-\kappa,r)=-H_r(\lambda,\kappa,r).
\end{equation}

The goal of this section is Theorem~\ref{thm:quadratic-fit}: at each
positive width, \((H,H_r)\) gives global coordinates on the model
parameter plane. We first fit the midpoint by varying the slope at
fixed curvature. Keeping the midpoint and width fixed, we then vary
the curvature and show that the model width derivative increases
through \((-1,1)\). The final subsection expresses the transport
of a general smooth profile in these fitted coordinates.

\subsection{Reference characteristics and curvature comparison}

It is useful to specify a quadratic arc by its slope at the peak.
For \(\kappa,z\in\mathbb R\), take the entire profile
\(zu+\kappa u^2/2\) with peak height zero.  As the base moves downwards
from that peak, let
\[
 j_\kappa(s,z),\qquad \Theta_\kappa(s,z),\qquad
 \mathcal L_\kappa(s,z)
\]
denote its base slope, its angle \(\operatorname{arctanh}(M/s)\), and
\(\operatorname{arctanh}v\), respectively, at characteristic width
\(s>0\).  Here \(v=M_r\) is the width derivative at fixed base height
of the given profile, evaluated on the characteristic.  Translating the
height coordinate and subtracting
the profile value at the translated origin shows that these relative
data describe every quadratic characteristic with peak slope \(z\)
and curvature \(\kappa\).

\begin{lemma}\label{fit:references}
The reference functions are smooth in \((\kappa,s,z)\) for \(s>0\) and
are defined at every finite positive width.  They satisfy
\begin{equation}\label{eq:fit-reference-system}
 \partial_s j_\kappa=-s\kappa\operatorname{sech}^2\Theta_\kappa,
 \qquad
 \partial_s\Theta_\kappa
   =j_\kappa-\frac{\sinh(2\Theta_\kappa)}s,
 \qquad
 \partial_s\mathcal L_\kappa
   =\frac{\sinh(2\Theta_\kappa)}s,
\end{equation}
with limiting values \(j_\kappa(0,z)=z\) and
\(\Theta_\kappa(0,z)=\mathcal L_\kappa(0,z)=0\).  Moreover,
\begin{equation}\label{eq:fit-reference-bounds}
 |j_\kappa(s,z)-z|\le\frac{|\kappa|s^2}{2},\qquad
 \partial_z j_\kappa(s,z)>0,\qquad
 \partial_z\Theta_\kappa(s,z)>0.
\end{equation}
For each fixed \(s>0\) and \(\kappa\), the map
\(z\mapsto\Theta_\kappa(s,z)\) is a bijection onto \(\mathbb R\).
Finally,
\begin{equation}\label{eq:fit-L-z}
 \partial_z\mathcal L_\kappa(s,z)
  =\int_0^s
    \frac{2\cosh(2\Theta_\kappa(q,z))}{q}
       \partial_z\Theta_\kappa(q,z)\,dq>0.
\end{equation}
\end{lemma}

\begin{proof}
The entire-profile construction in Lemma~\ref{lem:arc-data} gives an
arc at every lower base height.  Its width tends to zero at the peak
and tends to infinity as the base tends to minus infinity, since
\(s^2\ge 2(t-h)\).  Its strictly negative base derivative gives the
smooth inversion from base height to any \(s>0\), including smooth
dependence on the quadratic coefficients.  Thus no finite-width
existence claim is being inferred from the singular equations at
\(s=0\).  Equation~\eqref{eq:arc-characteristics} gives
Equation~\eqref{eq:fit-reference-system} and, by integration of the
slope equation, the first bound in Equation~\eqref{eq:fit-reference-bounds}.

For \(\kappa\ne0\), consider instead the single profile
\(\kappa u^2/2\).  Let \(\bar h(t,s)\) be its base height when the peak
height is \(t\) and the width is \(s\).  Its peak slope is \(z=\kappa t\),
and its base slope is \(\kappa\bar h(t,s)\).  Therefore
\[
 \partial_z j_\kappa(s,z)=\partial_t\bar h(t,s)>0
\]
by Lemma~\ref{lem:arc-data}.  This calculation is valid for either
sign of \(\kappa\).  For \(\kappa=0\), one has \(j_0(s,z)=z\).

The smooth scaling in Lemma~\ref{lem:arc-scaling} gives, locally
uniformly in the finite parameters,
\begin{equation}\label{eq:fit-reference-small-width}
 \Theta_\kappa(s,z)=\frac{z}{3}s+O(s^2),\qquad
 \partial_z\Theta_\kappa(s,z)=\frac{s}{3}+O(s^2).
\end{equation}
Differentiating the angle equation at positive width gives
\[
 \partial_s(\partial_z\Theta_\kappa)
  =\partial_z j_\kappa
     -\frac{2\cosh(2\Theta_\kappa)}s\partial_z\Theta_\kappa.
\]
The initial sign supplied by Equation~\eqref{eq:fit-reference-small-width}
and the positive forcing prove \(\partial_z\Theta_\kappa>0\) by the
scalar integrating-factor formula.

To verify the range, fix \(r>0\).  For \(0<s\le r\),
\[
 j_\kappa(s,z)\ge z-\frac{|\kappa|r^2}{2}.
\]
For sufficiently large positive \(z\), the angle is positive near
zero and cannot cross zero downwards, since its derivative there
would be positive.  Given \(N>0\), on \([r/2,r]\) and whenever
\(0\le\Theta_\kappa\le N\),
\[
 \partial_s\Theta_\kappa
 \ge z-\frac{|\kappa|r^2}{2}-\frac{2\sinh(2N)}r.
\]
For large enough \(z\), this lower bound forces a crossing of \(N\)
within that interval if the angle has not already exceeded \(N\);
it also prohibits a downward crossing of \(N\).  Thus
\(\Theta_\kappa(r,z)\to+\infty\) as \(z\to+\infty\).
Applying the same argument to \((-j,-\Theta,-\kappa)\) proves the
negative limit as \(z\to-\infty\).  Strict monotonicity and smoothness
give the asserted bijection.

Integrating the last equation in Equation~\eqref{eq:fit-reference-system}
gives
\[
 \mathcal L_\kappa(s,z)
   =\int_0^s\frac{\sinh(2\Theta_\kappa(q,z))}{q}\,dq.
\]
Equation~\eqref{eq:fit-reference-small-width}, including its locally
uniform parameter derivative, bounds the differentiated integrand
near zero.  Differentiation under this integral is therefore valid
and yields Equation~\eqref{eq:fit-L-z}.  Its integrand is strictly
positive at every positive width.
\end{proof}

For \(s>0\), let \(z=z(\kappa,s,\theta)\) be the inverse supplied by
Lemma~\ref{fit:references}, and define the reference slope field
\begin{equation}\label{eq:fit-slope-field}
 J_\kappa(s,\theta)
   =j_\kappa\bigl(s,z(\kappa,s,\theta)\bigr).
\end{equation}
It is smooth, and
\begin{equation}\label{eq:fit-field-positive}
 (J_\kappa)_\theta
   =\frac{\partial_z j_\kappa}{\partial_z\Theta_\kappa}>0.
\end{equation}
The following comparison will also be used for nonquadratic profiles.

\begin{lemma}[Curvature comparison]\label{lem:curvature-comparison}
Let a smooth profile \(\phi\) have an arc with peak height \(t\), and
parameterize that arc from its peak to a final width \(r>0\) by
\(s\in(0,r]\).  Assume \(\phi\) is smooth on a neighborhood of the
traversed height interval.  Write \(h(s)\), \(\theta(s)\), and \(L(s)\)
for its characteristic data, and set \(k=\phi'\).  For any fixed
\(\kappa\in\mathbb R\), define
\[
 E(s)=k(h(s))-J_\kappa(s,\theta(s)),\qquad
 \theta(s)=\Theta_\kappa(s,z(s)).
\]
Then \(z(s)\to k(t)\), \(E(s)\to0\) as \(s\downarrow0\), and
\begin{align}
 E'(s)+(J_\kappa)_\theta(s,\theta(s))E(s)
  &=-s\bigl(\phi''(h(s))-\kappa\bigr)
       \operatorname{sech}^2\theta(s),
       \label{eq:curvature-error}\\
 z'(s)&=\frac{E(s)}{\partial_z\Theta_\kappa(s,z(s))}.
       \label{eq:label-derivative}
\end{align}
Put \(a(s)=(J_\kappa)_\theta(s,\theta(s))>0\) and let \(f(s)\)
denote the right-hand side of Equation~\eqref{eq:curvature-error}.
The comparison from zero is given by the convergent formula
\begin{equation}\label{eq:fit-error-integral}
 E(s)=\int_0^s
       \exp\left(-\int_q^s a(\xi)\,d\xi\right)f(q)\,dq.
\end{equation}
In particular, if \(\phi''(h(s))\ge\kappa\) for \(0<s\le r\), then
\(E\le0\), \(z\) is nonincreasing, and
\[
 \theta(s)\ge\Theta_\kappa(s,z(r))\quad(0<s<r),\qquad
 L(r)\ge\mathcal L_\kappa(r,z(r)).
\]
If the curvature difference is not identically zero on \((0,r)\),
the last inequality is strict, as are all the preceding angle
inequalities for \(0<s<r\), and \(E(r)<0\).
All signs reverse under the assumption \(\phi''(h(s))\le\kappa\).
\end{lemma}

\begin{proof}
Along every reference characteristic, the definition of the field and
Equation~\eqref{eq:fit-reference-system} give
\[
 (J_\kappa)_s+(J_\kappa)_\theta
       \left(J_\kappa-\frac{\sinh(2\theta)}s\right)
   =-s\kappa\operatorname{sech}^2\theta.
\]
Subtracting this identity from the evolution of the actual slope in
Equation~\eqref{eq:arc-characteristics} proves
Equation~\eqref{eq:curvature-error}.  Differentiating
\(\theta(s)=\Theta_\kappa(s,z(s))\) and subtracting the two angle
equations gives Equation~\eqref{eq:label-derivative}.

The small-width expansions give \(\theta(s)/s\to k(t)/3\).
For each \(\varepsilon>0\), the reference angles with labels
\(k(t)-\varepsilon\) and \(k(t)+\varepsilon\) bracket this actual
angle at all sufficiently small positive widths.  The strict
monotonicity in Lemma~\ref{fit:references} therefore implies
\(z(s)\to k(t)\).  The bound
\(\lvert j_\kappa(s,z(s))-z(s)\rvert\le |\kappa|s^2/2\) then proves
\(E(s)\to0\).

For \(0<\varepsilon<s\), variation of constants gives
\[
 E(s)=E(\varepsilon)
       \exp\left(-\int_\varepsilon^s a\right)
      +\int_\varepsilon^s
       \exp\left(-\int_q^s a\right)f(q)\,dq.
\]
The first term tends to zero because \(a>0\) and
\(E(\varepsilon)\to0\).  Also \(f(q)=O(q)\) near the smooth peak,
and the exponential is between zero and one for \(0<q\le s\).
Dominated convergence proves Equation~\eqref{eq:fit-error-integral},
without any assumption that \(a\) is integrable at zero.

If the curvature difference is nonnegative, then \(f\le0\), so
\(E\le0\) and \(z'\le0\).  This proves the weak angle inequality
by Lemma~\ref{fit:references}.  If the difference is positive
somewhere in \((0,r)\), continuity gives a subinterval on which
\(f<0\).  Equation~\eqref{eq:fit-error-integral} makes \(E\) strictly
negative from the end of that subinterval through \(r\).  Thus
\(z(s)>z(r)\) for every \(s<r\), giving strict angle inequalities.
The integral formula for \(L\) in Equation~\eqref{eq:arc-L-integral}
then gives the stated strict inequality at the final width.
Negating the inequalities proves the other case.
\end{proof}

\subsection{Midpoint fitting and the positive conditional derivative}

For a prescribed \(r>0\), \(|M|<r\), and curvature \(\kappa\),
Lemma~\ref{fit:references} gives a unique peak label satisfying
\(\Theta_\kappa(r,z)=\operatorname{arctanh}(M/r)\).
Translate that reference arc to base height zero and subtract its
profile value at the base.  The resulting quadratic has base slope
\(\lambda=j_\kappa(r,z)\), curvature \(\kappa\), and midpoint \(M\).
Conversely, each base-zero quadratic arc has a finite peak and
arises in this way.  Since \(P>0\), midpoint fitting is unique.
We denote its slope by
\begin{equation}\label{eq:fit-midpoint-slope}
 \lambda_*(r,M,\kappa),\qquad
 H(\lambda_*(r,M,\kappa),\kappa,r)=M.
\end{equation}
The implicit-function theorem and uniqueness make \(\lambda_*\)
jointly smooth throughout \(r>0\), \(|M|<r\), \(\kappa\in\mathbb R\).
Define the conditional derivative and its angle by
\begin{equation}\label{eq:fit-conditional-data}
 \widehat v(r,M,\kappa)
  =H_r(\lambda_*(r,M,\kappa),\kappa,r),\qquad
 \widehat L(r,M,\kappa)=\operatorname{arctanh}\widehat v(r,M,\kappa).
\end{equation}
Both functions are smooth, by Equation~\eqref{eq:fit-open-data}.

\begin{lemma}\label{fit:conditional-derivative}
For every \(r>0\), \(|M|<r\), and \(\kappa\in\mathbb R\),
\begin{equation}\label{eq:fit-conditional-positive}
 \partial_\kappa\widehat L(r,M,\kappa)>0,
 \qquad
 \partial_\kappa\widehat v(r,M,\kappa)=\frac GP>0,
\end{equation}
where the model derivatives on the right are evaluated at the
midpoint fit.  In particular, \(G>0\) for every model arc.
\end{lemma}

\begin{proof}
Fix \(r,M,\kappa\).  For \(\delta>0\), let \(\theta_\delta(s)\)
be the characteristic of the midpoint-fitted model of curvature
\(\kappa+\delta\); let \(\theta_0\) be the one of curvature \(\kappa\).
Their final angles are equal.  Lemma~\ref{lem:curvature-comparison}
gives
\begin{equation}\label{eq:fit-interior-order}
 \theta_\delta(s)>\theta_0(s)\qquad(0<s<r).
\end{equation}
It follows already that the conditional \(\widehat L\) is strictly
increasing.  We prove a quantitative first-order bound to obtain
the asserted strict derivative.

Use the field \(J_\kappa\) to compare the curvature-\(\kappa+\delta\)
characteristic, and write its error as \(E_\delta\).  Then
\[
 E_\delta'+a_\delta E_\delta
    =-\delta s\operatorname{sech}^2\theta_\delta,
 \qquad
 a_\delta(s)=(J_\kappa)_\theta(s,\theta_\delta(s))>0,
 \qquad E_\delta\le0.
\]
Choose \(0<a<b<c<r\).  The already established smooth midpoint
fit and smooth finite-parameter arc dependence imply that, for all
sufficiently small \(\delta\ge0\), there are constants \(A,q_0>0\),
independent of \(\delta\), such that on \([a,r]\),
\[
 a_\delta(s)\le A,\qquad
 s\operatorname{sech}^2\theta_\delta(s)\ge q_0.
\]
Variation of constants starting at \(a\), using
\(E_\delta(a)\le0\), consequently gives
\begin{equation}\label{eq:fit-error-quantitative}
 E_\delta(s)\le-c_0\delta\quad(b\le s\le r),
 \qquad c_0=q_0(b-a)e^{-Ar}>0.
\end{equation}
This estimate needs no bound uniform in \(\delta\) near \(s=0\).
The sign at \(a\) follows from the from-zero comparison separately
for every \(\delta>0\).

Let \(z_\delta(s)\) be the reference label defined by
\(\theta_\delta(s)=\Theta_\kappa(s,z_\delta(s))\), and let \(z_0\)
be the constant label of the curvature-\(\kappa\) arc.  Final angle
matching gives \(z_\delta(r)=z_0\).  Smoothness of the inverse
reference-label map puts all these labels, for \(b\le s\le r\) and
small \(\delta\ge0\), in a common compact set.  Enlarge the set to
include the intervals between \(z_0\) and \(z_\delta(s)\).
Lemma~\ref{fit:references} then gives constants \(m_0,B_0>0\) with
\[
 m_0\le\partial_z\Theta_\kappa(s,z)\le B_0
\]
throughout that set.  Equations~\eqref{eq:label-derivative}
and~\eqref{eq:fit-error-quantitative} yield
\[
 z_\delta'(s)\le-\frac{c_0\delta}{B_0}\quad(b\le s\le r).
\]
Integrating backwards from the common final label gives, for
\(b\le s\le c\),
\[
 z_\delta(s)-z_0\ge\frac{c_0\delta(r-c)}{B_0},\qquad
 \theta_\delta(s)-\theta_0(s)
    \ge\frac{m_0c_0\delta(r-c)}{B_0}.
\]
Since the derivative of \(\sinh(2\theta)\) is at least \(2\),
the integral formulas for the final \(L\)-values and
Equation~\eqref{eq:fit-interior-order} now imply
\begin{equation}\label{eq:fit-L-quantitative}
 \widehat L(r,M,\kappa+\delta)-\widehat L(r,M,\kappa)
 \ge \frac{2m_0c_0(r-c)}{B_0}\log\!\frac cb\;\delta.
\end{equation}
Thus ordinary smoothness of the conditional function at the final
positive width proves \(\partial_\kappa\widehat L>0\).
No differentiation of the improper \(L\)-integral is used here.

Finally, implicit differentiation of
Equation~\eqref{eq:fit-midpoint-slope} gives
\(\partial_\kappa\lambda_*=-Q/P\), and hence
\[
 \partial_\kappa\widehat v
     =S-R\frac QP=\frac GP
     =(1-\widehat v^2)\partial_\kappa\widehat L>0.
\]
Since every model arc is its own midpoint fit, this proves \(G>0\)
everywhere.  The argument used \(P>0\), but did not assume the sign
or nonvanishing of \(G\).
\end{proof}

\subsection{The full range of the conditional fit}

\begin{lemma}\label{fit:properness}
For fixed \(r>0\) and \(|M|<r\),
\begin{equation}\label{eq:fit-properness}
 \lim_{\kappa\to+\infty}\widehat L(r,M,\kappa)=+\infty,
 \qquad
 \lim_{\kappa\to-\infty}\widehat L(r,M,\kappa)=-\infty.
\end{equation}
\end{lemma}

\begin{proof}
Suppose the first limit fails.  Strict monotonicity then makes
\(\widehat L(r,M,\kappa)\) bounded above as \(\kappa\to+\infty\).
Choose positive curvatures \(\kappa_n\uparrow+\infty\), and denote
the matched-final-angle characteristics by \(\theta_n(s)\),
\(j_n(s)\), and their final \(L\)-values by \(L_n\).
Lemma~\ref{lem:curvature-comparison} gives pointwise increase of
\(\theta_n(s)\) for \(0<s<r\).  Put
\[
 f_n(s)=\frac{\sinh(2\theta_n(s))}{s}.
\]
Each \(f_n\) is integrable at zero by its smooth-peak expansion,
and
\[
 f_n-f_1\ge0,\qquad
 \int_0^r(f_n-f_1)\,ds=L_n-L_1\le C
\]
for some finite \(C\).  Monotone convergence shows that the limiting
\(f_n\), and therefore the limiting \(\theta_n\), are finite almost
everywhere in \((0,r)\).

Choose four such widths \(0<a<b<c<d<r\).  The angle equation gives
\begin{align*}
 A_n:=\int_a^b j_n(s)\,ds
    &=\theta_n(b)-\theta_n(a)+\int_a^b f_n(s)\,ds,\\
 B_n:=\int_c^d j_n(s)\,ds
    &=\theta_n(d)-\theta_n(c)+\int_c^d f_n(s)\,ds.
\end{align*}
Both sequences are bounded above and below: the four endpoint
angle sequences are bounded, and for every subinterval \(I\),
\[
 0\le\int_I(f_n-f_1)\,ds\le C.
\]
Choose constants \(A_*,B_*\) with \(A_n\le A_*\) and
\(B_n\ge B_*\).  Because \(\kappa_n>0\), the slope \(j_n\) is
decreasing, so
\begin{equation}\label{eq:fit-outer-averages}
 j_n(b)\le\frac{A_n}{b-a}\le\frac{A_*}{b-a},\qquad
 j_n(c)\ge\frac{B_n}{d-c}\ge\frac{B_*}{d-c}.
\end{equation}
Thus \(j_n(b)-j_n(c)\) has a uniform upper bound.

On the other hand, its exact evolution is
\begin{equation}\label{eq:fit-middle-drop}
 j_n(b)-j_n(c)
    =\kappa_n\int_b^c s\operatorname{sech}^2\theta_n(s)\,ds.
\end{equation}
The integrand is bounded by \(s\) and converges almost everywhere
to a strictly positive function, because the limiting angle is
finite almost everywhere.  Dominated convergence gives a strictly
positive limiting integral in Equation~\eqref{eq:fit-middle-drop}.
Its right-hand side therefore tends to infinity, contradicting
Equation~\eqref{eq:fit-outer-averages}.  This proof uses no uniform
angle bound on the middle interval.

For the second limit, Equation~\eqref{eq:fit-reflection} and uniqueness
of midpoint fitting give
\[
 \lambda_*(r,-M,-\kappa)=-\lambda_*(r,M,\kappa),\qquad
 \widehat L(r,M,\kappa)=-\widehat L(r,-M,-\kappa).
\]
The positive-curvature result holds for every fixed midpoint,
including \(-M\).  It therefore proves the negative-curvature limit
at the original fixed midpoint \(M\), as required.
\end{proof}

\subsection{The remaining derivative signs and the global inverse}

\begin{lemma}\label{fit:positive-derivatives}
The model derivatives \(R\) and \(S\) are strictly positive for every
\((\lambda,\kappa,r)\in\mathbb R^2\times(0,\infty)\).
The width identity is
\begin{equation}\label{eq:fit-width-identity}
 \kappa P+\lambda=\frac{rH_r+H}{r^2-H^2}.
\end{equation}
With \(L_{\mathrm{mod}}(\lambda,\kappa,r)
=\operatorname{arctanh}H_r(\lambda,\kappa,r)\), differentiating
Equation~\eqref{eq:fit-width-identity} with respect to \(r\)
at fixed \((\lambda,\kappa)\) gives
\begin{equation}\label{eq:fit-R-identity}
 \kappa R
 =\frac{r(1-H_r^2)}{r^2-H^2}
   \left((L_{\mathrm{mod}})_r-\frac{2H}{r^2-H^2}\right).
\end{equation}
\end{lemma}

\begin{proof}
At base height \(h\), translating the quadratic and subtracting its
base value changes its slope to \(\lambda+\kappa h\) and leaves its
curvature unchanged.  Its midpoint data are therefore
\(H(\lambda+\kappa h,\kappa,r)\).  At \(h=0\), the first transport
identity in Equation~\eqref{eq:arc-transport} becomes
\[
 \kappa P-\frac r{r^2-H^2}H_r
     =\frac H{r^2-H^2}-\lambda,
\]
which is Equation~\eqref{eq:fit-width-identity}.
For completeness, put \(D_0=r^2-H^2\) and \(v=H_r\).
Differentiating at fixed \(\lambda,\kappa\) gives
\begin{align*}
 \kappa R
 &=\frac{rv_r+2v}{D_0}
      -\frac{2(r-Hv)(rv+H)}{D_0^2}\\
 &=\frac{r(1-v^2)}{D_0}
      \left(\frac{v_r}{1-v^2}-\frac{2H}{D_0}\right),
\end{align*}
proving Equation~\eqref{eq:fit-R-identity}.

Let \(t=t(\lambda,\kappa,r)>0\) be the peak of this base-zero model
arc.  Lemma~\ref{lem:arc-data} gives \(t_r>0\) at fixed model
parameters.  Its peak slope is \(z=\lambda+\kappa t\), so
\[
 L_{\mathrm{mod}}(\lambda,\kappa,r)
       =\mathcal L_\kappa(r,\lambda+\kappa t(\lambda,\kappa,r)).
\]
The chain rule and Equation~\eqref{eq:fit-reference-system} yield
\begin{equation}\label{eq:fit-base-width-L}
 (L_{\mathrm{mod}})_r
    =\frac{2H}{r^2-H^2}
       +\kappa t_r\partial_z\mathcal L_\kappa(r,z).
\end{equation}
For \(\kappa\ne0\), substitution into
Equation~\eqref{eq:fit-R-identity} and cancellation of \(\kappa\) give
\[
 R=\frac{r(1-H_r^2)}{r^2-H^2}
       t_r\partial_z\mathcal L_\kappa(r,z)>0
\]
by Lemma~\ref{fit:references}.  This cancellation is valid for
either sign of \(\kappa\).  For \(\kappa=0\), the peak slope is
\(z=\lambda\), so direct parameter differentiation gives
\[
 R=\partial_\lambda H_r
      =(1-H_r^2)\partial_z\mathcal L_0(r,\lambda)>0.
\]
Finally, \(S=(G+QR)/P>0\) by
Equation~\eqref{eq:fit-PQ-positive} and
Lemma~\ref{fit:conditional-derivative}.
\end{proof}

\begin{theorem}[Global quadratic fit]\label{thm:quadratic-fit}
For every \(r>0\) and every pair \((M,v)\) satisfying
\(|M|<r\) and \(|v|<1\), there is a unique finite pair
\((\lambda,\kappa)\in\mathbb R^2\) such that
\[
 H(\lambda,\kappa,r)=M,\qquad H_r(\lambda,\kappa,r)=v.
\]
The fitted parameters depend jointly smoothly on \((M,v,r)\) in
this open domain.  Equivalently, the smooth map
\begin{equation}\label{eq:fit-joint-map}
 (\lambda,\kappa,r)\longmapsto
       (H(\lambda,\kappa,r),H_r(\lambda,\kappa,r),r)
\end{equation}
is a diffeomorphism from \(\mathbb R^2\times(0,\infty)\) onto
\(\{(M,v,r):r>0,\ |M|<r,\ |v|<1\}\).
At every model arc, \(P,Q,R,S,G>0\).
In particular, fitted parameters converge to a finite pair whenever
their data converge to an interior point of this domain, even when
the width varies.
\end{theorem}

\begin{proof}
Fix \((M,v,r)\) in the stated domain.  For each curvature there is
exactly one slope fit by Equation~\eqref{eq:fit-midpoint-slope}.
Lemmas~\ref{fit:conditional-derivative} and~\ref{fit:properness}
show that its conditional \(\widehat L\) is a strictly increasing
map of \(\mathbb R\) onto \(\mathbb R\).  There is therefore exactly
one curvature for which \(\widehat L=\operatorname{arctanh}v\),
and then exactly one slope.  The sign assertions have been proved
in Equation~\eqref{eq:fit-PQ-positive} and
Lemmas~\ref{fit:conditional-derivative} and~\ref{fit:positive-derivatives}.

The Jacobian of Equation~\eqref{eq:fit-joint-map} has last row
\((0,0,1)\), and its determinant is \(G>0\).  The inverse-function
theorem gives a smooth local inverse at every point.  These inverses
agree wherever their domains overlap, by the global uniqueness
just proved, so the inverse is jointly smooth on the entire target
domain.  Its continuity gives the final convergence statement.
\end{proof}

\subsection{Transport of an actual profile in fitted coordinates}

We now express the motion of a general arc in the model coordinates.
Let $\phi\in C^\infty((0,\infty),\mathbb R)$, set $k=\phi'$,
and let $M(h,r)$ and $v(h,r)=M_r(h,r)$ be its endpoint data from
Section~\ref{sec:arcs}. The fit records both the midpoint and its
response to a change in width; their two transport equations determine
the evolution of the fitted slope and curvature.
Theorem~\ref{thm:quadratic-fit} gives unique smooth functions $\lambda(h,r)$ and $\kappa(h,r)$ such that
\begin{equation}\label{eq:fit-identities}
 M(h,r)=H(\lambda(h,r),\kappa(h,r),r),\qquad
 v(h,r)=H_r(\lambda(h,r),\kappa(h,r),r),
 \qquad h,r>0.
\end{equation}
In the following identities, $P,Q,R,S,G$ denote the quadratic-model
quantities evaluated at this fit.  In particular,
\[
 P,Q,R,S,G>0.
\]
Recall the notation
\[
 \alpha(h,r)=\frac{r}{r^2-M(h,r)^2},\qquad
 D=\partial_h-\alpha\partial_r.
\]
Along an arc with its peak fixed, $D$ is differentiation with respect
to the base height.  If the same characteristic is parametrized by its
outward-increasing width $s$, then differentiation along it is
$-D/\alpha$.

\begin{lemma}[Transport of the fit]\label{lem:fit-transport}
The fitted parameters satisfy
\begin{equation}\label{eq:fit-transport}
 \begin{aligned}
 P\lambda_r+Q\kappa_r&=0,\\
 D\kappa&=-\frac{R}{G}\bigl(\lambda-k(h)\bigr),\\
 D\lambda-\kappa&=\frac{S}{G}\bigl(\lambda-k(h)\bigr).
 \end{aligned}
\end{equation}
\end{lemma}

\begin{proof}
Differentiating the first identity in Equation~\eqref{eq:fit-identities}
at fixed $h$ and using the second one gives
$P\lambda_r+Q\kappa_r=0$.

For a quadratic profile, increasing its base height changes the slope
at rate $\kappa$, leaves its curvature fixed, and changes the width at
rate $-\alpha$.  Applying Equation~\eqref{eq:arc-transport} to that
profile, at the fitted data, therefore gives
\[
 \begin{aligned}
 \kappa P-\alpha H_r
   &=\frac{M}{r^2-M^2}-\lambda,\\
 \kappa R-\alpha H_{rr}
   &=-\frac{2Mr(1-v^2)}{(r^2-M^2)^2}.
 \end{aligned}
\]
For the actual profile, the chain rule gives
\[
 \begin{aligned}
 DM&=P D\lambda+Q D\kappa-\alpha H_r,\\
 Dv&=R D\lambda+S D\kappa-\alpha H_{rr}.
 \end{aligned}
\]
The expression for $Dv$ in Equation~\eqref{eq:arc-transport} depends
only on $M,v,r$, so it is identical for the two matched sets of data.
Subtracting the preceding identities consequently yields
\[
 \begin{aligned}
 P(D\lambda-\kappa)+Q D\kappa&=\lambda-k(h),\\
 R(D\lambda-\kappa)+S D\kappa&=0.
 \end{aligned}
\]
Their determinant is $G=PS-QR>0$.  Solving this system proves
Equation~\eqref{eq:fit-transport}.

\end{proof}
